\chapter{Introduction to Waves}
When we talk about waves, we refer to so many phenomena in nature. Waves can be in the form of liquids, like when you throw a pebble into a pond and create a recurring effect of water. Waves can be ocean waves, controlled by the moon and the tides. Waves can be sound waves, bouncing back and forth travelling from a speaker. Waves can even my magnetic or electric, which we call electromagnetic waves, creating charge and magnetic fields. Light waves creates the large range of frequency of light we see with the human eye. All of these waves, however, have something in common: \emph{they all share the same wave properties.}

We define waves as follows: 

\begin{quote}
A \textbf{wave} is a disturbance that transfers energy from one place to another without permanently transferring matter.
\end{quote}

Another way to put it: 
\begin{quote}
    A \textbf{wave} is a vibration of a system out of a restful state.
\end{quote}
For example, when a water wave moves across a pond, the water itself does not travel across the pond. Instead, the water particles oscillate around their equilibrium positions while the wave carries energy outward. 

In the case of a spring, which we will analyze deeper, energy (by the form of stretching or pushing) is transferred throughout the medium that is the spring.

\section{Basic Properties of Waves}

All waves can be described using several important properties. You will definitely see an overlap between trigonometric functions and wave properties. Often, we can describe waves that are generated with these trigonometric functions we have previously talked about.

There are essentially 5 terms we use to describe the properties of waves: \textbf{equilibrium}, \textbf{amplitude}, \textbf{wavelength}, \textbf{period}, and \textbf{frequency}.
\subsection*{Equilibrium}
An \textbf{equilibrium} state of a wave is what it would look no oscillation to be present. In other words, it is a state of rest for the wave. If you have a body of water that a pebble is thrown into, the equilibrium state is height or state of the water the instant before the pebble is thrown in. Same for a rope, a equilibrium state of a rope represents the rope before any motion is put into it.

\subsection*{Amplitude}

When a wave is created, it moves creates motion away from its \emph{equilibrium} position. 

The \textbf{amplitude} of a wave is the maximum displacement of the medium from its equilibrium position. In simpler terms, it measures how ``large'' the wave is.

The distance between the highest positive amplitude and the equilibrium position (called the \newterm{peak} or \newterm{crest}) and the lowest positive amplitude and the equilibrium position (called the \newterm{valley} or \newterm{trough}) are the same distance.

For a water wave, the amplitude is the height from the resting water level to the crest of the wave. For a sound wave, greater amplitude corresponds to louder sound. In light waves, larger amplitude corresponds to greater brightness. 

\subsection*{Wavelength}

The \textbf{wavelength}, usually represented by the symbol $\lambda$ (lambda), is the distance between two identical points on consecutive waves, such as crest-to-crest or trough-to-trough.

Wavelength is usually measured in meters (or centimeters or micrometers).

\subsection*{Frequency}

The \textbf{frequency} of a wave is the number of complete waves that pass a point each second. Frequency is measured in hertz (Hz), where

\[
1 \text{ Hz} = 1 \text{ wave per second}.
\]

A wave with a higher frequency oscillates more rapidly. For sound waves, higher frequency corresponds to higher pitch. For light waves, frequency determines the color of the light.

\subsection*{Period}

The \textbf{period} of a wave is the time required for one complete oscillation. The period and frequency are related by

where:
\begin{itemize}
    \item $T$ is the period (seconds),
    \item $f$ is the frequency (hertz).
\end{itemize}

\subsection*{Wave Speed}

Waves travel at a certain speed depending on the medium through which they move. Wave speed is related to frequency and wavelength by
$$v=f\lambda$$
where:
\begin{itemize}
    \item $v$ is wave speed,
    \item $f$ is frequency,
    \item $\lambda$ is wavelength.
\end{itemize}

ADD: exercise in labelling
% Considering adding formula for waves from sin/cos but would require parametric earlier on in the sequence
\section{Two main types of waves}

In waves around us, we can classify them under two categories: \textbf{longitudinal} and \textbf{transverse}. This classification depends on how particles are effected as the wave ``passes''. 

\subsection*{Longitudinal}
For \newterm{longitudinal} waves, the direction of wave (or the distrubance in particles) is \emph{parallel} to the direction of the vibration. The most common example of this a \emph{sound wave}. A speaker playing a sound pushes one sound particle, pushing another, and then another, ultimately continuing until the particles reach your ears.
%FIXME diagram
There is an amazing visual of this using a slinky in your digital resources!
\subsection*{Transverse}
In contrast, \newterm{transverse} waves move in a direction (or the distrubance in particles) is \newterm{perpendicular} to the wave's direction. Light and Electromagnetic waves and light waves are examples of transverse waves. The energy moves left to right (or horizontally) but the ends of the wave slide up and down vertically.
%FIXME diagram

Again, see your digital resources for a visual of this using a slinky.
\section{Compression and Rarefaction}

Sound waves travel by creating disturbances in a medium such as air. When an object vibrates, it pushes nearby air particles together in some regions and spreads them farther apart in others. These alternating regions move outward through the medium as the sound wave travels.

Regions where particles are crowded closer together are called \newterm{compressions}. In these regions, the pressure and density of the air are higher. Regions where particles are spread farther apart are called \newterm{rarefactions}\footnote{Rarefaction regions and phases are commonly called decompression regions.}. In these regions, the pressure and density are lower.

As the wave passes, the particles themselves do not travel long distances with the wave. Instead, they oscillate back and forth around their equilibrium positions while the disturbance moves through the medium.

\section{Creating a Wave in Python}

ADD ME